% =========================
% 5_evaluation.tex --- PLACEHOLDER. No evaluation has been run against the completed encoder.
% =========================
\section{Evaluation}
\label{sec:experiments}
\label{sec:methodology}
\label{sec:datasets}
\label{sec:eval_protocol}
\label{sec:label_efficiency}
\label{sec:pending_measurements}

\pendingsection{%
No recognition result exists for the system this paper specifies. The Phase-A encoder is complete
and its memory and train-only resolvability artifacts are built against it, but the admissibility
gate has not been fitted to those artifacts, so no enrollment curve, control comparison, or gate
telemetry is attributable to it. The corpus this section will evaluate on is described in
Appendix~\ref{app:corpus}.
\medskip

\textbf{This section will contain:}
\begin{itemize}\itemsep1pt
  \item the enrollment protocol --- support and query drawn from disjoint recorded executions,
        nested support prefixes at $k\in\{0,1,2,4,8\}$, a fixed candidate inventory per cell, and
        macro-F1 as the primary metric;
  \item development enrollment curves on MotionSense, RealWorld and Shoaib, reported against all
        five controls: ungated retrieval, support removal, cyclic support-label shuffling,
        prototypes, and ridge heads on the same frozen features;
  \item \emph{cross-configuration} cells, in which the enrolled configuration differs from the
        queried one. These are a precondition rather than an extension: a two-sided admissibility
        score degenerates to a per-candidate constant when both sides read the same sensor, so a
        same-configuration-only protocol cannot measure the mechanism (Section~\ref{sec:discussion});
  \item held-out concept, dataset, and body-region studies for the gate, with its gate-collapse,
        provenance-capture, and latent-to-measurement guards;
  \item ablations --- native-grid versus common-rate, sensor-folded versus channel-level,
        descriptor-on, acquisition-bias-on, observability-off, and signed low-frequency;
  \item data-matched and released-checkpoint baselines, including the short fine-tuning baseline
        the wearable-foundation-model literature requires; and
  \item after the development criteria are met, one evaluation of the sealed seven-dataset test
        roster with confidence intervals and seen/unseen candidate breakdowns.
\end{itemize}
\medskip

We claim no on-device latency, Core~ML, iOS, or unknown-class result: the corresponding
implementation and measurement artifacts do not exist in the repository.}
